\documentclass[aspectratio=169,11pt]{beamer}

\usepackage[utf8]{inputenc}
\usepackage[T1]{fontenc}
\usepackage[portuguese]{babel}
\usepackage{graphicx}
\usepackage{booktabs}
\usepackage{xcolor}

\usetheme{Madrid}
\usecolortheme{beaver}

% Hide navigation symbols
\setbeamertemplate{navigation symbols}{}

\title[Lynceus: Validação Científica]{Lynceus v2.0: Isolamento de Contribuição e Validação de Paridade em Motores eBPF}
\author[Herkenhoff, L.]{Leonardo Herkenhoff}
\institute[PPComp/IFES]{Mestrado Profissional em Computação Aplicada (PPComp) -- IFES}
\date{\today}

\begin{document}

\begin{frame}
    \titlepage
\end{frame}

\begin{frame}{O Problema Científico e a Crítica Acadêmica}
    \textbf{Crítica Recorrente na Revisão (O Desafio)}
    \begin{itemize}
        \item Extratores clássicos afirmam ser superiores, porém misturam ganhos de \textit{Performance} com ganhos de \textit{Riqueza de Features}.
        \item \textbf{A Dúvida:} O Lynceus (eBPF) é mais rápido porque é executado no Kernel, ou porque processa os fluxos de forma matematicamente distinta?
        \item \textbf{A Exigência de Isolamento:} Para provar ganho de velocidade irrefutável, é necessário avaliar a ferramenta sob uma \textbf{base metodológica justa} (mesmas características dos concorrentes).
    \end{itemize}
\end{frame}

\begin{frame}{A Metodologia de Validação: O Teorema da Paridade}
    \textbf{Resposta à Orientação:} Restrição Dimensional
    \vspace{0.3cm}
    
    Limitamos o motor Lynceus a extrair \textbf{apenas} os exatos mesmos atributos dos concorrentes, forçando um "campo de provas plano" (Level Playing Field).
    \vspace{0.3cm}
    
    \begin{itemize}
        \item \textbf{Paridade NFX (NetFeatureXtract):} Restrição absoluta para apenas 12 atributos volumétricos puros da Camada 3 (L3).
        \item \textbf{Paridade RustiFlow:} Expansão vetorial para 159 atributos limitados à Camada 4 (L4 TCP/UDP comportamental).
        \item \textbf{NetFlowLyzer (Base):} Extração plena original da nossa arquitetura com 495 atributos irrestritos (L2 a L7).
    \end{itemize}
\end{frame}

\begin{frame}{O Mecanismo da Velocidade (Fundamentação)}
    \textbf{Por que o Lynceus atinge uma vazão superior?}
    \vspace{0.3cm}
    
    Extratores legados (como o NFX) centralizam dados em um \textit{RingBuffer} global, induzindo extrema contenção (\textit{Lock Contention}) em sistemas Multi-Core, resultando em quedas por \textit{Out of Memory}.
    
    \vspace{0.3cm}
    \textbf{O Diferencial Arquitetural (\textit{Lockless}):}
    \begin{itemize}
        \item Alocação de vetores independentes por núcleo físico (\texttt{ARRAY\_OF\_MAPS}).
        \item O gancho XDP utiliza \texttt{bpf\_get\_smp\_processor\_id()} para despachar as métricas de pacote para processos de \textit{User-space} atrelados (\textit{Thread Affinity}).
        \item \textbf{Resultado:} Paralelismo intrínseco, sem necessidade de \textit{mutexes} ou bloqueios temporais. Escalonamento linear.
    \end{itemize}
\end{frame}

\begin{frame}{Contribuição 1: Performance em Base Justa (Extração)}
    Quando forçado a operar no mesmo espaço dimensional estreito de seus pares, o Lynceus prova a superioridade do modelo de extração Lockless no espaço do Kernel.
    
    \begin{table}[]
        \centering
        \resizebox{\textwidth}{!}{
        \begin{tabular}{lrrrr}
        \toprule
        \textbf{Ferramenta (Vetor Global)} & \textbf{Pacotes Proc.} & \textbf{Vazão Máx. (PPS)} & \textbf{RAM (MB)} & \textbf{Tolerância a Faltas} \\
        \midrule
        \textbf{Lynceus Base (495 Feat.)} & 171.263.883 & 125.314 & 12.166,60 & Sobreviveu (2.0M fluxos) \\
        \textbf{Lynceus NFX Parity (12 Feat.)} & 68.740.315 & \textbf{182.701} & \textit{Leve} & Sobreviveu (2.0M fluxos) \\
        NFX Original (User-space) & \textit{Abortado} & \textit{Abortado} & 2,60 & \alert{Falhou no 1º PCAP} \\
        \textbf{Lynceus Rust Parity (159 Feat.)} & 68.274.893 & \textbf{173.601} & 9.140,50 & Sobreviveu (2.0M fluxos) \\
        RustiFlow Original & 31.716.588 & 42.665 & 3.435,60 & Sobreviveu (Fragmentou) \\
        \bottomrule
        \end{tabular}
        }
    \end{table}
    
    \textbf{Isolamento da Contribuição:} O bypass arquitetural comprova seu valor. A arquitetura eBPF do Lynceus triplica a vazão (PPS) do estado da arte em Rust, livre de \textit{OOM Killers}.
\end{frame}

\begin{frame}{O Mecanismo da Pureza (Fundamentação)}
    \textbf{Por que as ferramentas perdem precisão com os mesmos atributos?}
    \vspace{0.3cm}
    
    Extratores legados sofrem com latência de fila (\textit{Queue Delay}) na pilha de rede do SO em tráfego de alta volumetria.
    \vspace{0.3cm}
    
    \textbf{A Falsa Miopia de ML:}
    \begin{itemize}
        \item O atraso introduz "ruído" temporal. Atributos cruciais como \textit{Inter-Arrival Time} e \textit{Active/Idle Time} são severamente corrompidos.
        \item \textbf{A Solução Lynceus:} Ao inspecionar os pacotes no \textbf{gancho XDP} (no instante em que tocam a placa de rede), o \textit{timestamp} capturado possui precisão de \textit{nanosegundos}, com 100\% de fidelidade temporal original (zero-pollution).
    \end{itemize}
\end{frame}


\begin{frame}{Contribuição 2: Pureza Estatística (Sem Vazamentos Lógicos)}
    A coleta prematura no XDP resolve as distorções temporais.
    
    \begin{table}[]
        \centering
        \resizebox{0.85\textwidth}{!}{
        \begin{tabular}{lccc}
        \toprule
        \textbf{Vetor de Ataque (Dataset)} & \textbf{RustiFlow (Original) F1} & \textbf{Lynceus (Restrito a Rust) F1} \\
        \midrule
        \textbf{DrDoS\_DNS} & 1.0000 & 0.9998 \\
        \textbf{DrDoS\_LDAP} & \alert{0.9271} & \textbf{0.9998} \\
        \textbf{PCAPv6 (IPv6 Custom)} & \alert{0.7619} & \textbf{0.9940} \\
        \textbf{DrDoS\_MSSQL} & 0.9999 & 0.9996 \\
        \bottomrule
        \end{tabular}
        }
    \end{table}
    
    \textbf{Isolamento da Contribuição:} Extraindo as \textbf{exatas mesmas 159 características}, o Lynceus corrige os falsos-negativos drásticos do RustiFlow nos vetores sensíveis a \textit{timing} (LDAP e IPv6). A precisão na fonte é uma contribuição inegável.
\end{frame}

\begin{frame}{O Teorema do Deslocamento Dimensional (Bônus)}
    Como o Random Forest mantem \textbf{F1 > 0.99} mesmo sofrendo uma amputação de 495 para apenas 12 atributos L3 (Paridade NFX)?
    \vspace{0.2cm}

    \begin{table}[]
        \centering
        \resizebox{0.9\textwidth}{!}{
        \begin{tabular}{lc|lc}
        \toprule
        \textbf{Paridade RustiFlow (159 attr L4)} & \textbf{Peso Gini} & \textbf{Paridade NFX (12 attr L3)} & \textbf{Peso Gini} \\
        \midrule
        \textbf{ACK\_Cnt} (Estatística TCP) & 9,07\% & \textbf{BwdPacketsRate} (Vazão Volátil) & 18,99\% \\
        \textbf{SYN\_Cnt} (Estatística TCP) & 9,04\% & \textbf{PacketsCount} (Volume Fixo) & 17,99\% \\
        \textbf{Win\_Median} (Saturação Buff) & 8,52\% & \textbf{TotalBytes} (Volume Fixo) & 15,00\% \\
        \bottomrule
        \end{tabular}
        }
    \end{table}

    \vspace{0.2cm}
    O Modelo Matemático \textbf{desloca sua ancoragem}: Sem acesso à taxonomia de flag TCP, a entropia probabilística migra naturalmente para a \textbf{aceleração temporal} do fluxo. A resiliência está na pureza da amostragem temporal do Lynceus, independente da dimensão espacial.
\end{frame}

\begin{frame}{Conclusão e Defesa de Tese}
    \textbf{Resumo de Impacto Inequívoco (Para o SBSeg):}
    \vspace{0.3cm}
    
    \begin{enumerate}
        \item \textbf{Performance Lockless (Engenharia):} A alocação per-core via \texttt{ARRAY\_OF\_MAPS} viabiliza monitoramento wire-speed sem perdas, mitigando a contenção OOM observada em ferramentas consolidadas.
        \item \textbf{Fidelidade Semântica (Ciência de Dados):} A coleta antecipada (XDP hook) erradica distorções de \textit{Queue Delay}, devolvendo precisão absoluta (F1 > 0.99) a matrizes de aprendizado sensíveis ao tempo.
        \item \textbf{Capacidade Dimensional (O(1)):} Arquitetura preparada para classificar agressivamente 495 *features* com \textit{Boundary Checks} eficientes, provendo escalabilidade sem comprometer o fluxo estocástico.
    \end{enumerate}
\end{frame}

\end{document}
